\documentclass[11pt,a4paper]{article}

\usepackage[T1]{fontenc}
\usepackage[utf8]{inputenc}
\usepackage{lmodern}
\usepackage[margin=2.5cm]{geometry}
\usepackage{microtype}
\usepackage{enumitem}
\usepackage{xcolor}
\usepackage{hyperref}

\setlength{\parindent}{0pt}
\setlength{\parskip}{0.6em}
\setlist[itemize]{topsep=0.2em,itemsep=0.2em}

\definecolor{linkblue}{HTML}{1F4E79}
\hypersetup{
    colorlinks=true,
    linkcolor=linkblue,
    urlcolor=linkblue,
    pdftitle={VIP Weeks 3--4 Report},
    pdfauthor={Yorgo Azzi}
}

\begin{document}

\begin{center}
    {\Large\bfseries VIP Weeks 3--4 Report}\\[0.5em]
    {\large Yorgo Azzi}\\[0.25em]
    September 30, 2026
\end{center}

\section{Overview}

During these two weeks, I replaced the colored-ribbon experiment with a
two-hand interface based on Google MediaPipe. OpenCV still captures and
displays the webcam image, while MediaPipe identifies both hands and returns
21 landmarks for each one. These landmarks provide more detailed movement and
gesture information than color-blob tracking and do not require the performer
to wear a colored object.

The main objective was to use the left hand for continuous controls and the
right hand for harmony. The interface was connected to the existing realtime
Markov and SATB generation engine so movement can affect the music without
interrupting playback.

\section{Work Completed}

\subsection{MediaPipe Hand Interface}

I added a MediaPipe Hand Landmarker model and a new realtime camera interface.
OpenCV reads and mirrors the frames, MediaPipe detects the hands, and the
program draws the landmarks and current control values over the camera image.
Hand labels are swapped by default for the current mirrored-camera setup, with
an option to disable the swap when necessary.

The program processes MediaPipe frames asynchronously so that old camera
frames can be dropped instead of building up delay. The music generator runs
on a separate clock, allowing camera analysis and MIDI playback to continue at
the same time.

\subsection{Left-Hand Musical Controls}

The left hand controls two continuous musical parameters:

\begin{itemize}
    \item palm height controls volume, with higher positions producing louder
    music;
    \item leftward hand tilt controls rhythmic density, with greater tilt
    favoring shorter notes and a faster-feeling rhythm.
\end{itemize}

The tilt angle is smoothed before it reaches the generator. Density continues
to bias the learned duration Markov model rather than forcing fixed note
values, so the generated rhythm remains variable.

\subsection{Right-Hand Chord Control}

The right hand selects one of the seven scale degrees through finger
combinations. The patterns progress from the index finger for the first degree
to combinations involving additional fingers, while thumb-only selects the
sixth degree and thumb plus index selects the seventh.

The right-hand tilt determines chord quality:

\begin{itemize}
    \item leaning left selects a minor chord;
    \item an upright hand uses the diatonic quality from the current key and
    mode;
    \item leaning right selects a major chord.
\end{itemize}

The selected key still defines the seven chord roots, but tilt can override
their normal diatonic quality. A stable gesture is accepted after a short dwell
time and applied at the next beat boundary. Unrecognized gestures and brief
tracking losses preserve the last valid chord.

\subsection{Gesture-Reliability Improvements}

Initial tests showed that individual fingers sometimes changed state for one
camera frame and that a tucked thumb was often interpreted as open. To improve
accuracy, I added per-finger temporal filtering, joint-angle and straightness
checks, short dropout tolerance, and a palm-relative thumb-distance test. The
camera overlay now marks extended fingertips in green and folded fingertips in
red, making calibration problems easier to identify.

\section{Results}

The completed prototype now connects the full interaction path:

\begin{center}
Camera $\rightarrow$ OpenCV $\rightarrow$ MediaPipe landmarks
$\rightarrow$ movement controls $\rightarrow$ Markov/SATB engine
$\rightarrow$ MIDI
\end{center}

The main results were:

\begin{itemize}
    \item Both hands can be tracked simultaneously without using a colored
    ribbon.
    \item Volume, density, scale degree, and chord quality can be controlled
    while the music continues playing.
    \item All 21 combinations of seven degrees and three quality modes produced
    valid SATB sonorities in automated testing.
    \item Finger filtering and thumb-to-palm distance improved chord-gesture
    stability during camera use.
    \item All 23 automated tests for camera controls, gesture recognition,
    chord selection, rhythm behavior, and SATB generation passed.
\end{itemize}

\section{Challenges}

Hand tracking is more expressive than colored-blob tracking, but it introduces
new calibration problems. Fingers can overlap, a thumb can remain visually
straight while folded across the palm, and detection may vary with camera
angle and distance. Mirrored video can also reverse MediaPipe's left/right
labels. These issues were reduced through configurable thresholds, temporal
filtering, visual feedback, and default hand-label swapping, but physical
testing remains important for tuning the interface to the performer.

\section{Next Steps}

\begin{itemize}
    \item Look for additional musical parameters that could be controlled by
    movement while keeping each mapping understandable.
    \item Display the generated sheet music in real time as the performance is
    being produced.
    \item Create a complete architecture diagram showing the connections among
    OpenCV, MediaPipe, movement analysis, the Markov models, SATB generation,
    MIDI scheduling, and sound output.
\end{itemize}

\section{Resources}

\begin{itemize}
    \item Project repository:
    \url{https://github.com/YorgoAzzi/MarkovModel}
    \item MediaPipe Hand Landmarker documentation:
    \url{https://ai.google.dev/edge/mediapipe/solutions/vision/hand_landmarker}
    \item OpenCV documentation:
    \url{https://docs.opencv.org/}
\end{itemize}

\end{document}
